% !TEX root = ../main.tex
\chapter{Binary Functions}
\label{ch:binary_functions}

Binary functions are elementwise operations on two inputs.
\Lucid\ ships approximately 50 of them, in four categories:
arithmetic, comparison, logical, and bitwise. Each is an
instantiation of \code{BinaryKernel<Op>}
(\chref{ch:kernel_abstraction}~Section~\ref{sec:kernel_abstraction:binary})
and inherits the kernel template's broadcasting and type-promotion
machinery automatically.

The interesting questions for binary functions are not the
individual forward expressions --- those are obvious --- but rather
the rules that govern how the kernel handles inputs of mismatched
shape, mismatched dtype, or mismatched device. We treat each rule
in turn, then enumerate the ops.

\secref{sec:binary_functions:broadcast} describes broadcasting.
\secref{sec:binary_functions:promotion} describes type promotion.
\secref{sec:binary_functions:device_check} describes the
cross-device rule.
\secref{sec:binary_functions:arithmetic}--\secref{sec:binary_functions:bitwise}
enumerate the op categories.

\section{Broadcasting}
\label{sec:binary_functions:broadcast}

Broadcasting is the rule by which two tensors of different but
\emph{compatible} shapes are extended to a common shape without
copying the data. \Lucid\ implements the NumPy broadcasting rule,
which is identical to \refframework's.

\subsection{The rule}
\label{ssec:binary_functions:broadcast_rule}

Given two shapes \(s_1\) and \(s_2\):

\begin{enumerate}
  \item Pad the shorter shape with leading 1s until both have the
    same rank.
  \item For each dimension, the output dimension size is the
    maximum of the two inputs at that dimension.
  \item The broadcast is legal if and only if, at each dimension,
    the two input sizes are either equal or one of them is 1.
  \item A dimension of size 1 in an input is repeated to match
    the corresponding output dimension --- but the repetition is
    \emph{virtual}: the input tensor is exposed through a stride of
    0 along the repeated dimension, with no data duplication.
\end{enumerate}

A \code{LucidError} is raised if the shapes are incompatible. The
error message gives both shapes and the dimension that
disagrees.

\subsection{Implementation: stride-0 expansion}
\label{ssec:binary_functions:stride_zero}

The mechanism by which broadcasting is free is the stride-0
trick. A tensor of shape \((3, 1, 5)\) broadcast to \((3, 4, 5)\)
is represented internally as a tensor of shape \((3, 4, 5)\) whose
strides for the broadcast dimension are zero --- every read along
that dimension lands at the same memory location. The kernel does
not need to know broadcasting happened; it iterates over the
output shape with the input's (modified) strides and the right
values flow through.

The cost is a small bit of stride-handling logic inside the kernel,
plus a one-time shape computation. There is no per-element overhead.

\subsection{Backward through broadcasting}
\label{ssec:binary_functions:broadcast_back}

If the forward broadcast a tensor from \((3, 1, 5)\) to \((3, 4,
5)\), the backward gradient (of shape \((3, 4, 5)\)) must be summed
along the broadcast dimension to match the original \((3, 1, 5)\)
input. \code{BinaryKernel} records the broadcast pattern at
forward time and applies the reverse-sum at backward time
automatically.

This is the \emph{unbroadcast} step --- it lives in the kernel
template, not in each op, since the pattern is identical for every
binary op.

\section{Type Promotion}
\label{sec:binary_functions:promotion}

When the two inputs have different dtypes, the kernel promotes
both to a common dtype before invoking the backend. The promotion
rules are described in
\chref{ch:backend_dispatch}~Section~\ref{sec:backend_dispatch:dtype_promo};
the short version is NumPy's rules.

\subsection{Scalar-tensor promotion}
\label{ssec:binary_functions:scalar_promo}

A common case is a tensor combined with a Python scalar
(\code{x + 1.0}). The scalar is wrapped into a 0-d tensor of
the appropriate dtype (\code{float32} for a Python \code{float},
\code{int64} for a Python \code{int}), then promotion runs as for
two tensors. The Python wrapper performs this conversion before
dispatching.

A subtle behaviour: \code{x + 1} where \code{x} is \code{float32}
does not promote to \code{int64}; the scalar 1 is interpreted in
\code{x}'s dtype context. This matches NumPy and prevents an
\code{int} addition from accidentally upcasting an FP32 tensor.

\section{Cross-Device Check}
\label{sec:binary_functions:device_check}

If the two input tensors are on different devices, the kernel
raises \code{LucidError("binary op across devices: a is on cpu,
b is on metal")}. \Lucid\ does not implicitly move tensors across
devices; the user must call \code{.to(device)} explicitly.

The rule applies to all binary ops, including comparisons. A
\code{Tensor.\_\_eq\_\_} of a CPU tensor and a metal tensor is an
error, not a silent move.

\section{Arithmetic}
\label{sec:binary_functions:arithmetic}

The arithmetic ops:

\begin{itemize}
  \item \code{add(a, b) = a + b}; backward: gradient flows to both
    inputs unchanged (modulo unbroadcast).
  \item \code{sub(a, b) = a - b}; backward: \code{grad\_a =
    grad\_out}, \code{grad\_b = -grad\_out}.
  \item \code{mul(a, b) = a * b}; backward: \code{grad\_a =
    b * grad\_out}, \code{grad\_b = a * grad\_out}.
  \item \code{div(a, b) = a / b}; backward: \code{grad\_a =
    grad\_out / b}, \code{grad\_b = -a / b\textasciicircum{}2 * grad\_out}.
  \item \code{pow(a, b)} returns $a^b$; backward: $\partial_a = b
    \cdot a^{b-1} \cdot \text{grad\_out}$, $\partial_b = a^b \cdot
    \ln(a) \cdot \text{grad\_out}$.
  \item \code{floordiv(a, b) = \lfloor a / b \rfloor}; integer
    division. Non-differentiable.
  \item \code{mod(a, b) = a - b * floordiv(a, b)}; remainder.
    Non-differentiable.
  \item \code{atan2(y, x) = \arctan(y/x)} with quadrant; backward:
    \code{grad\_y = x / (x\textasciicircum{}2 + y\textasciicircum{}2) * grad\_out}, \code{grad\_x =
    -y / (x\textasciicircum{}2 + y\textasciicircum{}2) * grad\_out}.
  \item \code{hypot(a, b)} returns $\sqrt{a^2 + b^2}$; backward derived
    accordingly.
  \item \code{logaddexp(a, b)} returns $\ln(e^a + e^b)$;
    numerically stable via the max-shift trick.
  \item \code{maximum(a, b), minimum(a, b)}; backward: gradient
    routes to the larger / smaller input, undefined at equality
    (we use 0.5 split by convention).
\end{itemize}

\subsection{The pow special case}
\label{ssec:binary_functions:pow_special}

\code{pow}'s backward expression $b \cdot a^{b-1}$ is undefined
for \code{a = 0, b < 1} (division by zero) and for negative \code{a}
with non-integer \code{b} (complex result). \Lucid\ returns zero
gradient at the undefined cases and propagates NaN otherwise,
matching \refframework. Parity tests check both behaviours.

\section{Comparison}
\label{sec:binary_functions:comparison}

\begin{itemize}
  \item \code{eq(a, b)}: elementwise equality. Returns bool.
  \item \code{ne(a, b)}: elementwise inequality.
  \item \code{lt(a, b), le(a, b), gt(a, b), ge(a, b)}: ordering
    comparisons.
\end{itemize}

Comparisons are not differentiable (the derivative of a step
function is a Dirac delta, which we represent as zero). The
autograd engine returns zero gradient for the inputs.

\subsection{Equality on floating-point}
\label{ssec:binary_functions:fp_eq}

\code{eq} on floating-point tensors is exact equality, not
tolerance-based. For tolerance comparisons, \code{isclose(a, b,
atol, rtol)} is the unary primitive that uses
\(|a - b| \leq \text{atol} + \text{rtol} \cdot |b|\).

\section{Logical}
\label{sec:binary_functions:logical}

The logical operators work on boolean tensors (or treat numeric
tensors as boolean, with nonzero as true):

\begin{itemize}
  \item \code{logical\_and(a, b)}: elementwise AND.
  \item \code{logical\_or(a, b)}: elementwise OR.
  \item \code{logical\_xor(a, b)}: elementwise XOR.
\end{itemize}

The unary \code{logical\_not} is in the unary chapter and follows
the same convention.

All logical ops are non-differentiable.

\section{Bitwise}
\label{sec:binary_functions:bitwise}

The bitwise operators work on integer tensors:

\begin{itemize}
  \item \code{bitwise\_and, bitwise\_or, bitwise\_xor}: per-bit
    operations.
  \item \code{lshift(a, b), rshift(a, b)}: arithmetic shifts.
\end{itemize}

Bitwise ops require both inputs to be integer dtypes; a
\code{LucidError} is raised if a float is passed. Non-differentiable.

\section{Operator Overloading}
\label{sec:binary_functions:overload}

Every arithmetic binary op is exposed both as a free function
(\code{lucid.add(a, b)}) and through Python operator overloading
on \code{Tensor} (\code{a + b}). The dunder methods
(\code{\_\_add\_\_}, \code{\_\_radd\_\_}, \code{\_\_iadd\_\_}, etc.)
route through the same dispatcher.

The in-place dunders (\code{\_\_iadd\_\_}, \code{\_\_imul\_\_})
invoke in-place ops if the tensor is not part of an autograd graph
and out-of-place ops otherwise. The choice avoids the in-place
version corrupting saved tensors for backward.

\subsection{Scalar interaction}
\label{ssec:binary_functions:scalar_interact}

The \code{\_\_radd\_\_}-style reflected dunders handle
\code{scalar + tensor} the same as \code{tensor + scalar}. The
scalar is wrapped to a 0-d tensor of the tensor's dtype, then the
normal binary op runs.

This is the convention that makes \code{1 + lucid.zeros(3)} work.

\section{The Broadcasting Algorithm in Detail}
\label{sec:binary_functions:broadcast_algo}

The high-level broadcasting rule given above hides a non-trivial
algorithm. This section walks through the algorithm step by step,
because the implementation in
\code{lucid/\_C/kernel/BinaryKernel.cpp} is one of the most
performance-critical pieces of code in the engine and the
correctness depends on getting every step right.

\subsection{Shape compatibility check}
\label{ssec:binary_functions:compat_check}

Given input shapes $s_a = (a_0, a_1, \ldots, a_{m-1})$ and
$s_b = (b_0, b_1, \ldots, b_{n-1})$ with $m \leq n$ (after
right-aligning by padding $s_a$ with leading 1s to length $n$),
the output shape is the elementwise maximum:
\[
o_i = \max(a_i', b_i) \quad \text{where } a_i' = \begin{cases}
1 & i < n - m \\
a_{i - (n-m)} & \text{otherwise}.
\end{cases}
\]
The broadcast is legal if and only if at each $i$, either
$a_i' = b_i$ or one of them equals 1.

The check runs in $O(n)$ time --- one pass over the rank, no per-
element work.

\subsection{Stride translation}
\label{ssec:binary_functions:stride_xlate}

To make the binary kernel iterate over the output shape while
reading the two inputs through their original strides, each
input's strides are mapped to the output rank as follows:

\begin{itemize}
  \item For dimensions present in the input
    (\(i \geq n - m\) for the shorter input), use the input's
    original stride.
  \item For dimensions added by padding, use stride 0.
  \item For dimensions where the input has size 1 but the output
    is larger, use stride 0 (the stride-0 broadcast trick).
\end{itemize}

The kernel then iterates over the output's $\prod_i o_i$ positions,
computing the linear offset into each input as the dot product of
the iteration index with the translated stride vector. The
arithmetic is simple but the addressing must be computed per
element; on modern Apple silicon CPUs the address computation
generates one MAC per dimension per element, which is the cost
of generality.

\subsection{Broadcast unbroadcast on backward}
\label{ssec:binary_functions:unbroadcast_algo}

The backward of a broadcast binary op must reverse the broadcast:
the upstream gradient has the output shape, but each input expects
a gradient of its original shape. The transformation is a sum
along each axis where the input was expanded:

\[
\text{grad\_a} = \sum_{\text{broadcast dims of } a} \text{grad\_out}.
\]

The kernel records the broadcast pattern at forward time (which
axes were stride-0 for each input) and the backward applies the
corresponding reduction. The reduction is fused into a single
kernel rather than a sequence of per-axis sums, which saves
intermediate allocations.

\section{Worked Examples}
\label{sec:binary_functions:examples}

We work through three representative binary op invocations to make
the broadcasting and unbroadcast concrete.

\subsection{Example: bias addition}
\label{ssec:binary_functions:bias_example}

A linear layer's bias addition:

\begin{lstlisting}[style=lucid,language=LucidPython]
# x.shape = (B, T, D); bias.shape = (D,)
out = x + bias        # broadcasts bias to (B, T, D)
\end{lstlisting}

Forward: the bias is right-aligned to shape $(1, 1, D)$ and stride-
zero-expanded to $(B, T, D)$. The kernel iterates over $B \cdot T
\cdot D$ positions, reading bias through stride 0 along the first
two axes (every read of bias along those axes lands at the same
memory).

Backward: \code{grad\_out} has shape $(B, T, D)$; the gradient for
the bias is summed over the first two axes:
\[
\text{grad\_bias} = \sum_{b, t} \text{grad\_out}_{b, t, :}.
\]
The kernel performs this reduction in a single fused pass.

\subsection{Example: outer product via broadcast}
\label{ssec:binary_functions:outer_example}

An outer product written as a broadcast multiplication:

\begin{lstlisting}[style=lucid,language=LucidPython]
# u.shape = (M,); v.shape = (N,)
outer = u[:, None] * v[None, :]   # shape (M, N)
\end{lstlisting}

Forward: \code{u[:, None]} has shape $(M, 1)$ with strides
$(\text{stride}_u, 0)$. \code{v[None, :]} has shape $(1, N)$ with
strides $(0, \text{stride}_v)$. Broadcasting expands both to
$(M, N)$: \code{u} is read with stride $(\text{stride}_u, 0)$ ---
the second axis still uses stride 0, so each row reads the same
$u_i$. Similarly for \code{v}.

The element at position $(i, j)$ of the output is $u_i \cdot v_j$
--- the outer product.

Backward: gradient \code{grad\_out} of shape $(M, N)$ flows back.
For \code{u} (which was broadcast along axis 1), sum over axis 1:
\code{grad\_u = grad\_out.sum(dim=1)}. For \code{v}, sum over
axis 0.

\subsection{Example: incompatible shapes}
\label{ssec:binary_functions:incompat_example}

A typical user mistake:

\begin{lstlisting}[style=lucid,language=LucidPython]
# x.shape = (B, T, D); y.shape = (B, D)
out = x + y     # error: T != D and neither is 1
\end{lstlisting}

After right-alignment, \code{y} becomes $(1, B, D)$. The two
shapes are $(B, T, D)$ and $(1, B, D)$. At axis 1, $T \neq B$ and
neither is 1; the broadcast is rejected. The error message names
both shapes and the offending axis.

\section{Complexity and Performance}
\label{sec:binary_functions:complexity}

Binary kernels are $O(N)$ in the total element count of the output
tensor. The constant factor varies by op:

\begin{itemize}
  \item Pure arithmetic (\code{add}, \code{mul}): one
    multiply-add per element, plus address computation.
  \item Transcendental (\code{atan2}, \code{logaddexp}): one
    transcendental function call per element. On Apple silicon
    these go through vForce's \code{vvatan2f}-style primitives or
    \MLX's \code{atan2} primitive, both well-optimised.
  \item Bitwise: one instruction per element, often the cheapest
    case.
\end{itemize}

\subsection{Memory bandwidth bound}
\label{ssec:binary_functions:bw_bound}

For tensors larger than the SLC (typically several megabytes),
binary kernels are memory-bandwidth-bound: the kernel reads two
inputs, writes one output, achieving roughly $3 N \cdot \text{dtype
size}$ bytes of memory traffic. On an M3~Max with 400~GB/s of
memory bandwidth, the upper bound for FP32 binary kernels is
roughly $400 / 12 \approx 33$ billion elements per second.

Measured throughputs on representative workloads come within
80--90\,\% of this bound for the contiguous case. The shortfall is
due to a mix of address-computation overhead and the kernel
launch cost.

\subsection{The stride-0 efficiency}
\label{ssec:binary_functions:stride0_eff}

A broadcast that uses stride-0 is faster than the equivalent
materialised broadcast (\code{repeat}). The reason is cache
behaviour: a stride-0 read along an axis hits the same cache line
repeatedly, while a materialised broadcast reads distinct cache
lines. For a $(B, 1, D)$ tensor broadcast to $(B, T, D)$ with $T =
1024$ and $D = 768$, the stride-0 broadcast reads $B D$ floats
($B \cdot 3$\,KB) while the materialised version reads $B T D$
floats ($B \cdot 3$\,MB) --- a thousand-fold reduction in DRAM
traffic.

The user-visible consequence: prefer broadcast over \code{repeat}
when the broadcast is read-only.

\section{Summary and Forward References}
\label{sec:binary_functions:summary}

The binary functions are roughly 50 ops across arithmetic,
comparison, logical, and bitwise categories, all instantiations of
\code{BinaryKernel<Op>}. Broadcasting follows NumPy rules and is
implemented through stride-0 expansion (free) with reverse-sum
unbroadcasting on the backward (also handled in the kernel
template). Type promotion follows NumPy rules. Cross-device
operations raise rather than implicitly move.

The next chapter, \chref{ch:reductions}, treats reductions ---
operations that collapse one or more axes of a tensor.
\chref{ch:shape_view} treats shape and view operations.
\chref{ch:indexing} treats indexing and gather/scatter.

\begin{lucidnote}
For the user: binary ops are the most ergonomic part of the API
because Python operator overloading makes them transparent. The
expression \code{x + y * z} compiles into three op calls with the
correct broadcasting, promotion, and autograd-graph construction
all handled automatically.
\end{lucidnote}
